% 18. Источник: lecture_13_expectation.tex: 17–33, 46–62, 67–74, 96–113. Страницы лекционного PDF: 56–57.
\question{Математическое ожидание и его смысл}

Пусть $\Omega=\{\omega_1,\ldots,\omega_n\}$. Проведём $N$ опытов
с результатами $\omega^1,\ldots,\omega^N$ и обозначим через $N_i$
число появлений исхода $\omega_i$. Тогда среднее наблюдаемых значений равно
\[
    \frac1N\sum_{j=1}^{N}\xi(\omega^j)
    =\sum_{i=1}^{n}\xi(\omega_i)\frac{N_i}{N}.
\]
В вероятностной модели частотам $N_i/N$ соответствуют вероятности
$\Pr\{\omega_i\}$.

\begin{definition}[\topicmark{Конечная модель}Математическое ожидание в конечной модели]
    \textbf{Математическим ожиданием} случайной величины $\xi$ называется
    \[
        \EE\xi=\sum_{i=1}^{n}\xi(\omega_i)\Pr\{\omega_i\}.
    \]
    Используется также обозначение $M\xi$.
\end{definition}

% Тервер(Лекция)68.pdf, с. 2; Тервер(Лекция)69.pdf, с. 1.
В общем вероятностном пространстве конечное математическое ожидание записывается как
$\EE\xi=\int_\Omega \xi(\omega)\,\diff\Pr(\omega)$ при $\int_\Omega|\xi|\,\diff\Pr<\infty$.

\begin{theorem}[\topicmark{Вычисление по распределению}Вычисление по распределению]
    Если в конечной модели $\xi$ принимает различные значения
    $x_1,\ldots,x_k$ с вероятностями $p_i=\Pr\{\xi=x_i\}$, то
    \[
        \EE\xi=\sum_{i=1}^{k}x_i p_i.
    \]
\end{theorem}

\begin{Proof}
    Сгруппируем исходы, на которых $\xi$ принимает одинаковое значение:
    \[
        \sum_{\omega\in\Omega}\xi(\omega)\Pr\{\omega\}
        =\sum_{i=1}^{k}x_i
          \sum_{\omega:\,\xi(\omega)=x_i}\Pr\{\omega\}
        =\sum_{i=1}^{k}x_i p_i.
    \]
\end{Proof}

\begin{definition}[\topicmark{Дискретный случай}Дискретное распределение]
    Для распределения $\{(x_i,p_i)\}_i$ математическое ожидание равно
    \[
        \EE\xi=\sum_i x_i p_i,
        \qquad\text{если}\quad \sum_i |x_i|p_i<\infty.
    \]
    Таким образом, требуется абсолютная сходимость ряда.
\end{definition}

\begin{definition}[\topicmark{Абсолютно непрерывный случай}Абсолютно непрерывное распределение]
    Если $\xi$ имеет плотность $p_\xi$, то
    \[
        \EE\xi=\int_{-\infty}^{+\infty}x p_\xi(x)\,\diff x,
        \qquad\text{если}\quad
        \int_{-\infty}^{+\infty}|x|p_\xi(x)\,\diff x<\infty.
    \]
    Если интеграл записан как несобственный интеграл Римана,
    требуется его абсолютная сходимость.
\end{definition}

Для малого промежутка $\Delta_i$ около точки $x_i$, в которой плотность
непрерывна,
\[
    \Pr\{\xi\in\Delta_i\}\approx p_\xi(x_i)|\Delta_i|.
\]
Поэтому суммам $\sum_i x_i\Pr\{\xi\in\Delta_i\}$ соответствуют
интегральные суммы $\sum_i x_i p_\xi(x_i)|\Delta_i|$.

% Тервер(Лекция)68.pdf, с. 3.
\begin{example}[Математическое ожидание пуассоновского распределения]
Если $\xi\sim\operatorname{Poiss}(\lambda)$, то слагаемое при $k=0$ равно нулю. В остальных слагаемых сокращаем $k$ и полагаем $m=k-1$:
\[
 \EE\xi=e^{-\lambda}\sum_{k=1}^{\infty}k\,\frac{\lambda^k}{k!}
 =\lambda e^{-\lambda}\sum_{m=0}^{\infty}\frac{\lambda^m}{m!}
 =\lambda.
\]
\end{example}

% 19. Источник: lecture_13_expectation.tex: 125–191. Страницы лекционного PDF: 58–59.
\question{Свойства математического ожидания}

\begin{theorem}[\topicmark{Ожидание константы}Ожидание константы]
    Для $c\in\R$ имеем $\EE c=c$.
\end{theorem}

\begin{Proof}[Дискретный случай\editorialmark]
У постоянной случайной величины единственное значение $c$ имеет
вероятность $1$. Поэтому сумма по её значениям состоит из одного
слагаемого: $\EE c=c\cdot 1=c$.
\end{Proof}

\begin{theorem}[\topicmark{Линейность}Линейность]
Если существуют конечные $\EE\xi_1$ и $\EE\xi_2$, то для любых
    $a_1,a_2\in\R$ существует математическое ожидание их линейной комбинации и
    \[
        \EE(a_1\xi_1+a_2\xi_2)
        =a_1\EE\xi_1+a_2\EE\xi_2.
    \]
\end{theorem}

\begin{Proof}[Линейность в дискретном случае]
    Пусть $\xi_1$ принимает значения $x_i$, а $\xi_2$ --- значения $y_j$.
    Обозначим
    \[
        p_{ij}=\Pr\{\xi_1=x_i,\ \xi_2=y_j\},\qquad
        p_i=\Pr\{\xi_1=x_i\},\quad q_j=\Pr\{\xi_2=y_j\}.
    \]
    По счётной аддитивности вероятности
    \[
        \sum_j p_{ij}=p_i,\qquad \sum_i p_{ij}=q_j.
    \]
    Сначала проверим абсолютную сходимость:
    \[
        \begin{aligned}
            \sum_{i,j}|a_1x_i+a_2y_j|p_{ij}
            &\leq |a_1|\sum_i|x_i|\sum_jp_{ij}
                  +|a_2|\sum_j|y_j|\sum_ip_{ij}\\
            &=|a_1|\sum_i|x_i|p_i+|a_2|\sum_j|y_j|q_j<\infty.
        \end{aligned}
    \]
    Следовательно, суммы можно перегруппировать:
    \[
        \begin{aligned}
            \EE(a_1\xi_1+a_2\xi_2)
            &=\sum_{i,j}(a_1x_i+a_2y_j)p_{ij}\\
            &=a_1\sum_i x_i p_i+a_2\sum_j y_j q_j\\
            &=a_1\EE\xi_1+a_2\EE\xi_2.
        \end{aligned}
    \]
\end{Proof}

\begin{theorem}[\topicmark{Монотонность}Монотонность]
    Если $\xi\geq\eta$ почти наверное и оба математических ожидания конечны,
    то $\EE\xi\geq\EE\eta$.
\end{theorem}

\begin{Proof}
    Для $\zeta=\xi-\eta\geq0$ имеем $\EE\zeta\geq0$.
    По линейности $\EE\zeta=\EE\xi-\EE\eta$, откуда
    $\EE\xi\geq\EE\eta$.
\end{Proof}

\begin{theorem}[\topicmark{Модуль ожидания}Оценка модуля ожидания]
    Если $\EE\xi$ конечно, то
    \[
        |\EE\xi|\leq\EE|\xi|.
    \]
\end{theorem}

\begin{Proof}
    Ожидание $\EE|\xi|$ конечно по условию абсолютной интегрируемости.
    Из $-|\xi|\leq\xi\leq|\xi|$, монотонности и линейности следует
    \[
        -\EE|\xi|\leq\EE\xi\leq\EE|\xi|.
    \]
\end{Proof}

% 20. Источник: lecture_17_expectation.tex: 7–33, 37–61, 66–84. Страницы лекционного PDF: 60–61.
\question{Ожидание произведения. Неравенства Йенсена и Ляпунова}

\begin{theorem}[\topicmark{Ожидание произведения}Ожидание произведения независимых величин]
    Если $\xi$ и $\eta$ независимы и имеют конечные математические
    ожидания, то $\EE(\xi\eta)$ существует и
    \[
        \EE(\xi\eta)=\EE\xi\cdot\EE\eta.
    \]
\end{theorem}

\begin{Proof}[Дискретный случай]
    Пусть $\Pr\{\xi=x_i\}=p_i$ и $\Pr\{\eta=y_j\}=q_j$.
    По независимости
    \[
        \Pr\{\xi=x_i,\eta=y_j\}=p_iq_j.
    \]
    Абсолютная сходимость следует из равенства
    \[
        \sum_{i,j}|x_iy_j|p_iq_j
        =\left(\sum_i|x_i|p_i\right)
         \left(\sum_j|y_j|q_j\right)<\infty.
    \]
    Поэтому
    \[
        \EE\xi\cdot\EE\eta
        =\left(\sum_i x_ip_i\right)\left(\sum_j y_jq_j\right)
        =\sum_{i,j}x_iy_jp_iq_j=\EE(\xi\eta).
    \]
\end{Proof}

\begin{theorem}[\topicmark{Неравенство Йенсена}Неравенство Йенсена]
    Пусть $g\colon\R\to\R$ --- выпуклая вниз функция,
    а $\xi$ имеет конечное математическое ожидание. Тогда
    \[
        g(\EE\xi)\leq\EE g(\xi).
    \]
    Правая часть может быть равна $+\infty$.
\end{theorem}

\begin{Proof}
    Для любого $x_0\in\R$ существует число $\ell(x_0)$, такое что
    \[
        g(x)\geq g(x_0)+(x-x_0)\ell(x_0),\qquad x\in\R.
    \]
    Положим $x_0=\EE\xi$ и $x=\xi$:
    \[
        g(\xi)\geq g(\EE\xi)+(\xi-\EE\xi)\ell(\EE\xi).
    \]
    По монотонности и линейности математического ожидания
    \[
        \EE g(\xi)\geq g(\EE\xi)
        +\ell(\EE\xi)\underbrace{\EE(\xi-\EE\xi)}_{=\,0}
        =g(\EE\xi).
    \]
\end{Proof}

\begin{theorem}[\topicmark{Неравенство Ляпунова}Неравенство Ляпунова]
    Если $0<s<t$ и $\EE|\xi|^t<\infty$, то
    \[
        \bigl(\EE|\xi|^s\bigr)^{1/s}
        \leq\bigl(\EE|\xi|^t\bigr)^{1/t}.
    \]
\end{theorem}

\begin{Proof}
    Из $|\xi|^s\leq1+|\xi|^t$ следует $\EE|\xi|^s<\infty$.

    \smallskip
    Функция $g(x)=|x|^{t/s}$ выпукла, поскольку $t/s>1$.
    Применяя неравенство Йенсена к $\xi'=|\xi|^s$, получаем
    \[
        \bigl(\EE|\xi|^s\bigr)^{t/s}
        =g(\EE\xi')\leq\EE g(\xi')
        =\EE\bigl(|\xi|^s\bigr)^{t/s}=\EE|\xi|^t.
    \]
    Возводя обе части в степень $1/t$, получаем требуемое неравенство.
\end{Proof}

% 21. Источник: lecture_17_expectation.tex: 66–84, 88–106. Страницы лекционного PDF: 61.
\question{Моменты и связь между их порядками}

\begin{definition}[\topicmark{Определение момента}Определение]
    \textbf{Моментом порядка $k$} случайной величины $\xi$, где
    $k\in\N$, называется число
    \[
        \EE\xi^k,
    \]
    если это математическое ожидание конечно, то есть
    $\EE|\xi|^k<\infty$.
\end{definition}

\begin{consequence}[\topicmark{Существование моментов}Существование моментов]
    Из неравенства Ляпунова следуют утверждения:
    \begin{enumerate}
        \item Если существует конечный момент порядка $k$, то существуют
        все моменты меньших натуральных порядков $i<k$.
        \item Если конечного момента порядка $k$ нет, то нет и конечных
        моментов натуральных порядков $i>k$.
    \end{enumerate}
\end{consequence}

% Тервер(Лекция)69.pdf, с. 3–4: оба следствия записаны из одного неравенства.
\begin{Proof}[\topicmark{Связь порядков моментов}Связь порядков моментов]
    В доказательстве неравенства Ляпунова при $0<s<t$ используется оценка
    \[
        |\xi|^s\leq 1+|\xi|^t.
    \]
    Из $\EE|\xi|^t<\infty$ следует $\EE|\xi|^s<\infty$.

    \smallskip
    То же утверждение в форме контрапозиции: если $\EE|\xi|^s=\infty$,
    то для любого $t>s$ имеем $\EE|\xi|^t=\infty$.
\end{Proof}

% 22. Источник: lecture_13_expectation.tex: 208–251; lecture_15_covariance.tex: 31–84. Страницы лекционного PDF: 62–64.
\question{Дисперсия, среднеквадратическое отклонение и ковариация}

\begin{definition}[\topicmark{Дисперсия}Определение]
    \textbf{Дисперсией} случайной величины $\xi$ с конечным математическим
    ожиданием называется
    \[
        \DD\xi=\operatorname{Var}\xi
        =\EE(\xi-\EE\xi)^2.
    \]
    Конечная дисперсия существует, если $\EE\xi^2<\infty$.
\end{definition}

\begin{definition}[\topicmark{Среднеквадратическое отклонение}Среднеквадратическое отклонение]
    \textbf{Среднеквадратическим отклонением} случайной величины $\xi$
    называется число
    \[
        \sigma_\xi=\sqrt{\DD\xi}.
    \]
\end{definition}

\begin{theorem}[\topicmark{Формула для дисперсии}Формула для вычисления дисперсии]
    При $\EE\xi^2<\infty$ справедливо равенство
    \[
        \DD\xi=\EE\xi^2-(\EE\xi)^2.
    \]
\end{theorem}

\begin{Proof}
    По линейности математического ожидания
    \[
        \begin{aligned}
            \DD\xi
            &=\EE\bigl(\xi^2-2\xi\EE\xi+(\EE\xi)^2\bigr)\\
            &=\EE\xi^2-2(\EE\xi)^2+(\EE\xi)^2\\
            &=\EE\xi^2-(\EE\xi)^2.
        \end{aligned}
    \]
\end{Proof}
\begin{theorem}[\topicmark{Линейное преобразование}Дисперсия линейного преобразования]
    Если $\DD\xi<\infty$, то для любых $a,b\in\R$
    \[
        \DD(a\xi+b)=a^2\DD\xi.
    \]
\end{theorem}

\begin{Proof}[Центрирование\editorialmark]
Положим $\eta=a\xi+b$. По линейности $\EE\eta=a\EE\xi+b$,
поэтому $\eta-\EE\eta=a(\xi-\EE\xi)$. По определению дисперсии
\[
  \DD\eta=\EE\bigl[(\eta-\EE\eta)^2\bigr]
  =a^2\EE\bigl[(\xi-\EE\xi)^2\bigr]=a^2\DD\xi.
\]
\end{Proof}

\begin{definition}[\topicmark{Ковариация}Определение]
    Для случайных величин $\xi$ и $\eta$ с конечными вторыми моментами
    \textbf{ковариацией} называется число
    \[
        \operatorname{cov}(\xi,\eta)
        =\EE\bigl[(\xi-\EE\xi)(\eta-\EE\eta)\bigr].
    \]
\end{definition}

Далее предполагаем, что у рассматриваемых величин конечны вторые моменты.

\begin{theorem}[\topicmark{Свойства ковариации}Свойства ковариации]
\[
 \operatorname{cov}(\xi,\xi)=\DD\xi,\qquad
 \operatorname{cov}(\xi,\eta)=\EE(\xi\eta)-\EE\xi\cdot\EE\eta.
\]
\end{theorem}
\begin{Proof}[Вычисление по определению\editorialmark]
При $\eta=\xi$ под знаком ожидания в определении ковариации стоит
$(\xi-\EE\xi)^2$. Его ожидание равно $\DD\xi$.

\smallskip
Для второй формулы положим $m=\EE\xi$ и $n=\EE\eta$. Раскрывая скобки, получаем
\[
 \EE[(\xi-m)(\eta-n)]
 =\EE(\xi\eta)-n\EE\xi-m\EE\eta+mn=\EE(\xi\eta)-mn.
\]
\end{Proof}

\begin{theorem}[Билинейность ковариации]
Для $a,b\in\R$ справедливы равенства
\[
 \begin{aligned}
 \operatorname{cov}(a\xi+b\eta,\zeta)
 &=a\operatorname{cov}(\xi,\zeta)+b\operatorname{cov}(\eta,\zeta),\\
 \operatorname{cov}(\zeta,a\xi+b\eta)
 &=a\operatorname{cov}(\zeta,\xi)+b\operatorname{cov}(\zeta,\eta).
 \end{aligned}
\]
\end{theorem}
\begin{Proof}[Линейность по каждому аргументу\editorialmark]
Обозначим $U^\circ=U-\EE U$. По линейности ожидания
$(a\xi+b\eta)^\circ=a\xi^\circ+b\eta^\circ$. Поэтому
\[
 \operatorname{cov}(a\xi+b\eta,\zeta)
 =\EE[(a\xi^\circ+b\eta^\circ)\zeta^\circ]
 =a\EE(\xi^\circ\zeta^\circ)+b\EE(\eta^\circ\zeta^\circ).
\]
Это первое равенство. Второе следует из симметрии ковариации:
$\EE(U^\circ V^\circ)=\EE(V^\circ U^\circ)$.
\end{Proof}

\begin{theorem}[Независимость и ковариация]
Из независимости $\xi$ и $\eta$ следует $\operatorname{cov}(\xi,\eta)=0$.
Обратное утверждение неверно.
\end{theorem}
\begin{Proof}[Независимые величины\editorialmark]
По свойству ожидания произведения независимых величин
$\EE(\xi\eta)=\EE\xi\cdot\EE\eta$. Подстановка в формулу для ковариации даёт нуль.
\end{Proof}
\begin{example}[Зависимые величины с нулевой ковариацией\editorialmark]
Пусть $\xi$ принимает значения $-1,0,1$ с вероятностями $1/3$,
а $\eta=\xi^2$. По симметрии $\EE\xi=\EE\xi^3=0$.
Поскольку $\xi\eta=\xi^3$, имеем $\operatorname{cov}(\xi,\eta)=0$.

\smallskip
События $\{\xi=0\}$ и $\{\eta=0\}$ совпадают. Вероятность их пересечения
равна $1/3$, а произведение вероятностей равно $(1/3)^2=1/9$. Следовательно,
$\xi$ и $\eta$ зависимы.
\end{example}

\begin{theorem}[\topicmark{Дисперсия суммы}Дисперсия суммы]
    Если $\xi$ и $\eta$ имеют конечные вторые моменты, то
    \[
        \DD(\xi+\eta)=\DD\xi+2\operatorname{cov}(\xi,\eta)+\DD\eta.
    \]
    В частности, для независимых $\xi$ и $\eta$
    \[
        \DD(\xi+\eta)=\DD\xi+\DD\eta.
    \]
\end{theorem}

\begin{Proof}
    По линейности математического ожидания
    \[
        \begin{aligned}
            \DD(\xi+\eta)
            &=\EE\bigl[(\xi-\EE\xi)+(\eta-\EE\eta)\bigr]^2\\
            &=\EE(\xi-\EE\xi)^2
              +2\EE\bigl[(\xi-\EE\xi)(\eta-\EE\eta)\bigr]
              +\EE(\eta-\EE\eta)^2\\
            &=\DD\xi+2\operatorname{cov}(\xi,\eta)+\DD\eta.
        \end{aligned}
    \]
    При независимости ковариация равна нулю.
\end{Proof}

% 23. Источник: lecture_17_correlation.tex: 8–36. Страницы лекционного PDF: 65.
\question{Корреляция и её свойства}

\begin{definition}[\topicmark{Коэффициент корреляции}Определение]
    Пусть $0<\DD\xi<\infty$ и $0<\DD\eta<\infty$.
    \textbf{Коэффициентом корреляции} случайных величин $\xi$ и $\eta$
    называется число
    \[
        \operatorname{corr}(\xi,\eta)
        =\frac{\operatorname{cov}(\xi,\eta)}
        {\sqrt{\DD\xi\cdot\DD\eta}}.
    \]
\end{definition}

\begin{theorem}[\topicmark{Свойства корреляции}Свойства коэффициента корреляции]
При конечных положительных дисперсиях $|\operatorname{corr}(\xi,\eta)|\leq1$.
\end{theorem}

% ТВ 3.pdf, с. 1–2; Тервер(Лекция)69.pdf, с. 5–6 (формулировки свойств).
\begin{Proof}[\topicmark{Оценка модуля корреляции}Оценка модуля корреляции]
Центрируем и нормируем обе величины:
\[
 X=\frac{\xi-\EE\xi}{\sqrt{\DD\xi}},\qquad
 Y=\frac{\eta-\EE\eta}{\sqrt{\DD\eta}}.
\]
Тогда $\EE X^2=\EE Y^2=1$ и $\operatorname{corr}(\xi,\eta)=\EE(XY)$.
Из $(|X|-|Y|)^2\geq0$ получаем $2|XY|\leq X^2+Y^2$. Значит,
\[
 |\operatorname{corr}(\xi,\eta)|=|\EE(XY)|
 \leq\EE|XY|\leq \tfrac12(\EE X^2+\EE Y^2)=1.
\]
\end{Proof}
\begin{theorem}[Корреляция независимых величин]
Если $\xi$ и $\eta$ независимы и их дисперсии конечны и положительны,
то $\operatorname{corr}(\xi,\eta)=0$.
\end{theorem}
\begin{Proof}[Нулевая корреляция\editorialmark]
При независимости $\operatorname{cov}(\xi,\eta)=0$ по свойству из вопроса~22.
В определении корреляции знаменатель $\sigma_\xi\sigma_\eta$ положителен,
поэтому $\operatorname{corr}(\xi,\eta)=0/(\sigma_\xi\sigma_\eta)=0$.
\end{Proof}

\begin{theorem}[Критерии корреляции $+1$ и $-1$]
При тех же условиях на дисперсии положим $\rho=\operatorname{corr}(\xi,\eta)$.
\[
 \begin{aligned}
 \rho=1  &\ \Longleftrightarrow\ \exists\, a>0,\ b\in\R:\ \xi=a\eta+b\quad\text{почти наверное},\\
 \rho=-1 &\ \Longleftrightarrow\ \exists\, a<0,\ b\in\R:\ \xi=a\eta+b\quad\text{почти наверное}.
 \end{aligned}
\]
\end{theorem}

\begin{Proof}[Случай равенства в оценке]
При $|\operatorname{corr}(\xi,\eta)|=1$ в предыдущей цепочке все оценки
становятся равенствами. В частности, $\EE|XY|=1$, поэтому
\[
 \EE\bigl[(|X|-|Y|)^2\bigr]
 =\EE X^2+\EE Y^2-2\EE|XY|=0.
\]
Следовательно, $|X|=|Y|$ почти наверное.
\end{Proof}
\begin{Proof}[$\Rightarrow$ Знак линейной связи\editorialmark]
Пусть $\rho=s$, где $s\in\{-1,1\}$. Выше установлено, что
$|X|=|Y|$ почти наверное и $\EE|XY|=1$. При этом $\EE(XY)=s$.
Величина $|XY|-sXY$ неотрицательна, а её ожидание равно $1-s^2=0$.
Следовательно, $sXY=|XY|$ почти наверное.

\smallskip
При $s=1$ знаки $X$ и $Y$ совпадают там, где они ненулевые,
а при $s=-1$ они противоположны. Вместе с $|X|=|Y|$ это даёт
$X=sY$ почти наверное. Возвращаясь к исходным величинам, получаем
\[
 \xi=a\eta+b,\qquad
 a=s\,\frac{\sigma_\xi}{\sigma_\eta},\qquad b=\EE\xi-a\EE\eta.
\]
Оба стандартных отклонения положительны, так что знак $a$ равен $s$.
\end{Proof}
\begin{Proof}[$\Leftarrow$ От линейной связи к корреляции\editorialmark]
Пусть $\xi=a\eta+b$ почти наверное и $a\ne0$. Тогда
$\xi-\EE\xi=a(\eta-\EE\eta)$ и $\sigma_\xi=|a|\sigma_\eta$.
Поэтому $X=(a/|a|)Y$. Так как $\EE Y^2=1$, имеем
\[
 \rho=\EE(XY)=\frac{a}{|a|}\EE Y^2
 =\begin{cases}1,&a>0,\\-1,&a<0.\end{cases}
\]
\end{Proof}
